\section*{4. 循环移位问题}

给定一个按从小到大排序的数组，现将其循环右移若干位（如原数组[1,2,3,4,5]，循环右移2位后为[4,5,1,2,3]），设计算法计算循环右移的位数k（$k \ge 0$），要求时间复杂度尽可能优化。

\begin{itemize}
    \item \textbf{核心性质}：循环右移后的数组是"两段有序子数组"的拼接，第一段所有元素$\ge$第二段所有元素，循环右移位数k等于最小值元素的索引
    \item \textbf{分治查找}：类似折半搜索，比较A[mid]与A[high]，若A[mid]$\le$A[high]则右半有序，最小值在左半；否则最小值在右半
    \item \textbf{终止条件}：low == high时找到最小值，或当A[mid]小于左右邻居时mid即为最小值位置
    \item \textbf{边界处理}：数组未移位（仍完全有序）时，最小值索引为0，k=0
    \item \textbf{时间复杂度}：每次排除一半，O(log n)
\end{itemize}

\begin{lstlisting}[language=C++, style=cppstyle]
int findRotationCount(vector<int>& A) {
    int n = A.size();
    int low = 0, high = n - 1;
    if (A[low] <= A[high]) return 0;
    while (low <= high) {
        int mid = (low + high) / 2;
        int prev = (mid - 1 + n) % n;
        int next = (mid + 1) % n;
        if (A[mid] <= A[prev] && A[mid] <= A[next]) {
            return mid;}
        if (A[mid] <= A[high]) {
            high = mid - 1;
        } else {
            low = mid + 1;}}
    return 0;}
\end{lstlisting}

\begin{enumerate}
    \item 先判断数组是否根本没移位（首元素$\le$尾元素），是的话直接返回0
    \item 找到中点，看它是不是最小值：比左边和右边都小，那就是了
    \item 如果不是，判断哪半边是有序的：右半有序说明最小值在左半，左半有序说明最小值在右半
    \item 向无序的那半边继续找（用二分的方式）
    \item 找到最小值的位置就是循环右移的位数k
\end{enumerate}
